\section{Fixed finite data and the final contradiction}\label{sec:frobenius}

The scalar bound and the two-slot jets will now be tested against one
independent theorem. We state all its hypotheses before constructing
the remaining data. Its independent proof in
\cite{OpenAIAbundance2026} uses neither abundance nor logarithmic
subadditivity.

\begin{theorem}[Finite-data Frobenius incompatibility]
\label{thm:finite-data}
Let $W$ be a smooth connected projective complex variety of dimension
$n\ge2$. Fix rational Cartier divisors $L,H$, with $H$ ample, an
integral Cartier divisor $S$, an integer $q>0$, and real numbers
$r>1$, $\epsilon>0$. Suppose
\begin{gather*}
 H^n\le(2\epsilon)^n,\qquad
 K_WH^{n-1}\le 2H^n/r,\qquad
 (2L+qS)H^{n-1}\le4H^n,\\
 (14\epsilon)^n<1/4,\qquad80\epsilon<3/4.
\end{gather*}
The following three kinds of fixed data cannot all exist.
\begin{enumerate}[label=\textup{(\roman*)}]
\item A smooth integral complete-intersection flag ending in a curve
$C$, of successive divisor classes $l_iH$ for positive integers $l_i$,
with $\mu_{\min}(\Omega_W^1|_C)\ge0$.
\item On $Z_W=\PP(\cO(S)_1\oplus\cO(S)_2)\to W\times W$,
in the symmetric-power convention, put $M=L_1+L_2+q\xi$.
For some integer $k_0>0$ with $k_0L$ Cartier, finitely many sections of
$k_0M$ surject onto
$\cO_{Z_W,z_*}/\mathfrak m_{z_*}^{(2n+2)k_0+1}$ after
trivialization at a smooth torus point $z_*$.
\item A point $x\in W$, the blowup $\pi_x:\widehat W\to W$
with exceptional divisor $J_x$, and a curve class
$\gamma=f_*(A_1\cdots A_{n-1})$, where
$f:V\to\widehat W$ is a fixed birational morphism from a smooth
integral projective variety and the $A_i$ are ample Cartier divisors,
such that $J_x\gamma>0$ and
\[
 \pi_x^*(L+K_W/2+\lambda S)\gamma
 \le40\epsilon J_x\gamma\qquad(0\le\lambda\le q).
\]
It suffices to check the two endpoint inequalities.
\end{enumerate}
No nefness of $L,S,K_W$ is assumed, and $z_*$ need not lie over $x$.
All divisors, the flag, the point, the curve representative, and the
finite jet sections are fixed over $\C$ before the residual
characteristic varies.
\end{theorem}

The proof is \cite[Theorem~9.1]{OpenAIAbundance2026}.
We shall apply it with $n=4$, $S=S_0=\iota K$, and $r=r_t$.
The remaining work is to obtain its ample perturbation, flag, and
actual movable curve witness from the fourfold geometry.

\subsection{Fixing the polarization, flag, and movable test}

Choose a positive rational number $\epsilon$ sufficiently small that
\begin{equation}\label{eq:epsilon-choice}
 (14\epsilon)^n<\frac14,\qquad 80\epsilon<\frac34.
\end{equation}
Choose $q$ as in \Cref{prop:two-slot-jets}, then a sufficiently small
rational $t>0$ so that both jet estimates hold,
$r_t>q\iota+1$, and $L^n<(2\epsilon)^n$. Fix $t$ and $r_t$. Let
$W\to Y$ be a projective resolution, and continue to denote
pullbacks by $K,A,L$.  Terminality gives
\[
 K_W=K+E,\qquad E\ge0\text{ exceptional},\qquad S_0=\iota K\text{ Cartier}.
\]
Since $L$ is pulled back from an ample divisor on $Y$,
\[
 K_WL^{n-1}=KL^{n-1}\le L^n/r_t.
\]
Perturb $L$ by a sufficiently small ample rational class on $W$.
By continuity, the resulting ample rational divisor $H$ satisfies
\begin{equation}\label{eq:fixed-H}
 H^n\le(2\epsilon)^n,\qquad
 K_WH^{n-1}\le\frac{2H^n}{r_t},\qquad
 (2L+qS_0)H^{n-1}\le4H^n.
\end{equation}
For the last inequality, at $H=L$ use
$S_0L^{n-1}\le \iota L^n/r_t$ and $q\iota<r_t$.
We now fix $H$; all further geometric
data will be chosen once over $\C$ before reduction.

The divisor $K_W$ is pseudo-effective.  Generic semipositivity, in
the smooth empty-boundary case of
\cite[Theorem~2.1]{CampanaPaun2015}, gives nonnegative
$H^{n-1}$-slope to every positive-rank torsion-free quotient of
$\Omega_W^1$.  Apply the Mehta--Ramanathan restriction theorem to the
Harder--Narasimhan filtration and its factors
\cite[Theorem~6.1 and Remark~6.2]{MehtaRamanathan1982}.  We obtain a
fixed flag of smooth integral general very ample sections, of types
\[
 h_i=l_iH\quad(1\le i\le n-1),
\]
ending in a smooth curve $C$, such that
\begin{equation}\label{eq:fixed-curve-slope}
 \mu_{\min}(\Omega_W^1|_C)\ge0.
\end{equation}
The $l_i$ are sufficiently large divisible integers, chosen
successively.  The final general curve avoids the codimension-two
loci where the filtration factors fail to be locally free.  Set
$\Lambda=\prod_{i=1}^{n-1}l_i$.  Then
\begin{equation}\label{eq:flag-degrees}
 [C]=\Lambda H^{n-1},\qquad h_i\cdot C=l_i\Lambda H^n.
\end{equation}

Next choose a very general point $x\in W$ above the isomorphism locus
of $W\to Y$, and let
$\pi_x:\widehat W\to W$ be its blowup, with exceptional divisor $J_x$.
Put
\[
 D^+=2r_t(K+2tA)+2E.
\]
By \Cref{prop:scalar-orders}, every nonzero section of a sufficiently
divisible multiple $kD^+$ satisfies
\begin{equation}\label{eq:Dplus-order}
 \ord_x(s)\le32k\epsilon.
\end{equation}
Indeed, pushing down removes only the effective exceptional part,
and is an isomorphism near $x$.

\begin{lemma}\label[lemma]{lem:fixed-movable-test}
There is a strongly movable curve class $\gamma$ on $\widehat W$,
represented as the pushforward of an ample complete intersection on
a fixed smooth projective birational model, such that
\begin{equation}\label{eq:movable-test}
 (\pi_x^*D^+)\cdot\gamma
 <40\epsilon\,J_x\cdot\gamma,\qquad J_x\cdot\gamma>0.
\end{equation}
\end{lemma}

\begin{proof}
If $\pi_x^*D^+-40\epsilon J_x$ were pseudo-effective, then for any
rational $c$ with $32\epsilon<c<40\epsilon$ the divisor
$\pi_x^*D^+-cJ_x$ would be big: it is a positive convex combination
of the preceding pseudo-effective divisor and the big divisor
$\pi_x^*D^+$.  A sufficiently divisible effective multiple would
contradict \Cref{eq:Dplus-order}.  Thus
$\pi_x^*D^+-40\epsilon J_x$ is not pseudo-effective.

Movable-curve duality \cite[Theorem~2.2]{BDPP2013} gives a strongly
movable class with negative pairing against that divisor.  Since the
inequality is strict, we may take an actual pushforward of an ample
complete intersection on one smooth birational projective model,
approximating the ample classes rationally and clearing their
denominators if needed.  The big divisor $\pi_x^*D^+$ has positive
pairing with this nonzero class: write it as an ample class plus an
effective class and pull back to the chosen model.  The strict
inequality therefore implies $J_x\cdot\gamma>0$.
\end{proof}

We omit $\pi_x^*$ in subsequent intersection formulas involving
$\gamma$.  The concrete complete-intersection representative will
continue to test effective divisors after reduction.

\subsection{Endpoint comparison}

For $0\le\lambda\le q$, put $Q_\lambda=L+K_W/2+\lambda S_0$.
The actual divisor classes on the fixed resolution satisfy
\begin{equation}\label{eq:slot-domination}
 D^+-Q_\lambda
 =(r_t-1/2-\lambda\iota)K+3r_ttA+\tfrac32E.
\end{equation}
The first coefficient is positive because $r_t>q\iota+1$.
The pullbacks of $K,A$ are nef, and $E$ is effective. A strongly
movable curve pairs nonnegatively with each of these classes.
Consequently \Cref{lem:fixed-movable-test} gives
\[
 Q_\lambda\gamma\le D^+\gamma
 <40\epsilon J_x\gamma\qquad(0\le\lambda\le q).
\]
This verifies the curve inequalities in \Cref{thm:finite-data}.
It does not require a pseudo-effective cone to specialize to
positive characteristic: the curve has the fixed ample
complete-intersection representative required by that theorem.

Take $z_*,k_0$, and the finite sections supplied by
\Cref{cor:fixed-two-slot-jets} on this resolution. For clarity about
the finite choices in the comparison, choose also an integral ample
divisor $T_0$ so that each $T_0+aS_0$, for
$0\le a\le(k_0-1)q$, has fixed sections nonvanishing at the two
base coordinates of $z_*$. Serre generation permits these finitely
many choices. They supply the residual degrees when the common
proof puts $B_p=k_0\lfloor p/k_0\rfloor L+T_0$. The differences
$B_p-pL$ belong to a fixed finite list.

We have now verified all the hypotheses of \Cref{thm:finite-data}.
The choices were made in the order $n=4,\epsilon$; then the gap
$\delta$ and integer $q$; then $t,r_t$; then the resolution,
$H$, flag, point, and curve; and finally the finite jets and fillers.
Only after these choices does its proof let $p$ tend to infinity.

\subsection{The two orders of the determinant}

We recall the comparison in the proof of
\cite[Theorem~9.1]{OpenAIAbundance2026}, identifying the determinant
bundles to which our movable test applies. Spread the fixed data as
in that proof and work on a smooth reduction in sufficiently large
characteristic $p$. The same notation denotes the reduced divisors
and witnesses. For relative Frobenius $F:W\to W'$, let $S_0'$ be
the base-field twist of $S_0$ and put
\[
 \mathcal E=F_*\cO_W(B_p),\qquad
 s=\rk\mathcal E=p^n,\qquad R=s^2=p^8,\qquad
 U_a=H^0(W,\cO_W(B_p+paS_0)).
\]
The projection formula gives an evaluation map
\begin{equation}\label{eq:finite-evaluation}
 \Phi_p:\bigoplus_{\substack{a+b=q\\a,b\ge0}}
 U_a\otimes U_b\otimes\cO_{W'\times W'}(-aS'_{0,1}-bS'_{0,2})
 \longrightarrow\mathcal E\boxtimes\mathcal E.
\end{equation}
Products of the fixed jet sections, multiplied by the fixed fillers,
generate the local quotient by the $p$-th powers of the parameters
at $z_*$. Let $A_*$ be the tensor product of the two base
Frobenius-fiber algebras. If the torus coordinate $z$ has value
$z_0\ne0$ at $z_*$, this quotient is
$A_*[z]/(z^p-z_0^p)$. Project to the coefficient of $1$ in its
$A_*$-basis $1,z,\ldots,z^{p-1}$. Exactly the weights divisible
by $p$ survive, showing that the values of
\eqref{eq:finite-evaluation} span $A_*$. Its dimension is $R$,
so $\Phi_p$ has generic rank $R$.

On the diagonal, the columns instead restrict to global sections of
one line bundle on the Frobenius fiber product:
\[
 \Delta_F=W\times_{W'}W,\qquad
 \mathcal L_p=
 \cO_{\Delta_F}(B_{p,1}+B_{p,2}+pqS_{0,1}).
\]
Here $pS_{0,1}$ and $pS_{0,2}$ are isomorphic line bundles on
$\Delta_F$, so all the weight pairs $a+b=q$ give this same bundle.
The rank at every diagonal point is at most
$h^0(\Delta_F,\mathcal L_p)$. The diagonal-ideal filtration has
graded bundles
\[
 \mathcal G_j=\cO_W(2B_p+pqS_0)\otimes A_j(\Omega_W^1),
 \qquad 0\le j\le n(p-1),
\]
where $A_j$ is the degree-$j$ part of the symmetric algebra modulo
the $p$-th powers of its local generators. If $e_j=\rk A_j$, then
$\sum_j e_j=p^n$. The slope estimate and our one fixed flag give
the common bound proved in the companion:
\[
 h^0(W,\mathcal G_j)
 \le e_jp^n\bigl(7^nH^n+O(p^{-1})\bigr).
\]
The error is independent of $j$. Summing therefore bounds every
diagonal rank by
$R(7^nH^n+O(p^{-1}))<R/4$, since
$7^nH^n\le(14\epsilon)^n<1/4$.

Choose $R$ generically independent columns of $\Phi_p$, and let
$m_i$ be the sum of their weights in slot $i$. Their nonzero
determinant is a section
\[
 0\ne\sigma\in H^0(W'\times W',\mathcal Q_1\boxtimes\mathcal Q_2),
 \qquad
 \mathcal Q_i=(\det\mathcal E)^{\otimes s}
                 \otimes\cO_{W'}(m_iS_0').
\]
These are actual line bundles. Under the numerical identification
with the base-field twist, their first Chern classes satisfy
\begin{equation}\label{eq:finite-determinant-classes}
 \frac{c_1(\mathcal Q_i)}R
 =\frac{B_p}p+\frac{p-1}{2p}K_W+\lambda_iS_0
 =Q_{\lambda_i}+\frac{B_p-pL-K_W/2}{p},
 \qquad \lambda_i=\frac{m_i}{R}\in[0,q].
\end{equation}
The numerator of the last term belongs to a fixed finite list.
Thus the endpoint comparison bounds the intersection of each
normalized class with $\gamma$ by
$(40\epsilon+O(p^{-1}))J_x\gamma$, uniformly in the chosen columns.

This intersection bound controls sections on the reduction itself.
Put $x'=F(x)$. For any nonzero section of $\mathcal Q_i$, the strict
transform of its zero divisor on the blowup at $x'$ is effective and
pairs nonnegatively with the
twist of $\gamma$: that curve is still the pushforward of the fixed
ample complete intersection. Since $J_x\gamma>0$, every such section
has order at $x'$ at most $R(40\epsilon+O(p^{-1}))$.
Bases adapted to vanishing order in the two factors show that a
nonzero section of their external product has order at most the sum
of these bounds. Within each bidegree the leading tensors are linearly
independent, and different bidegrees cannot cancel. Consequently
\[
 \ord_{(x',x')}\sigma\le R(80\epsilon+O(p^{-1})).
\]
But the rank of $\Phi_p$ at $(x',x')$ is less than $R/4$.
Constant row operations make more than $3R/4$ rows of the selected
matrix vanish at that point, forcing
$\ord_{(x',x')}\sigma\ge3R/4$. The two orders contradict
$80\epsilon<3/4$ for large $p$. The generic-rank calculation used
$z_*$, whereas this diagonal comparison uses the independently
chosen point $x$.

Thus the terminal counterexample in \Cref{prop:reduction} cannot
exist. Canonical section spaces are preserved by its birational
construction, so this proves \Cref{thm:nonvanishing}.
